% Thesis abstract.  A fragment: emitted by SCRIPTS/build_tex.py after \maketitle
% (see --abstract), so it carries no preamble of its own.

\begin{abstract}
\noindent
Pump-and-treat remediation of a contaminated aquifer is a design problem under
uncertainty: well rates must keep the concentration at a compliance surface below
a limit with high probability, while the conductivity field that decides where
the plume travels is known only statistically. The problem was
posed to me as a doctoral thesis in Applied Mathematics at UC Davis twenty-five
years ago, when the simulation cost made it impossible to solve as written; this
thesis asks it again with the tools that did not exist then, and holds the
answers to a standard the earlier attempt could not afford. It solves that problem
for a two-dimensional benchmark ($200\times120$\,m, $160\times96$ cells, a
log-normal conductivity truncated at $1\,529$ Karhunen--Lo\`eve modes, nine
candidate wells, a $5\%$ chance constraint) and treats the question of what makes
a computed design believable as a first-class object of study.

Three layers are built and verified in turn. A one-dimensional
advection--dispersion sandbox with the Ogata--Banks solution as its reference
supports the learned surrogates: a physics-informed network reaches a median
full-grid relative error of $1.61\times10^{-2}$ over five seeds (spread
$1.17$--$8.14\times10^{-2}$) and a coordinate-augmented Fourier neural operator
$1.99\times10^{-2}$ at 250 epochs, $11\times$ better than a mean-field predictor.
A batched finite-volume Darcy/transport solver with a parity-tested GPU backend
supports the optimization: a conditional-value-at-risk penalty in a sample
average approximation, minimized by compass search over well rates in
$1.67$--$4.14$\,h of one rented GPU. Two runs differing only in floating-point
precision reach policies $1.2\%$ apart in cost on different well subsets
($J = 428\,592$ and $433\,945$), each sitting on the constraint.

The certification layer supplies both the strongest confirmation and the
strongest negative result. The accepted policy reproduces its cost to $0.02\%$ on
four hundred scenarios no run had seen, with reliability $0.978$ (Wilson interval
$[0.958,0.988]$) against a $0.94$ gate; a policy optimized on only sixteen
scenarios reports perfect reliability in-run and fails out of sample
($0.912$ $[0.881,0.936]$). Across three disjoint draws of four hundred, however,
the CVaR tail moves between $+0.616$, $+0.254$ and $+0.334$ --- spread smaller
than the standard error of any one value, because a $5\%$ tail of four hundred
draws averages twenty of the most variable realizations. Reliability transports;
the tail does not. The thesis therefore recommends certifying with reliability
and its interval, quoting tails with their standard error, and never ranking
policies by tail value at affordable sample sizes.

\vspace{0.6em}
\noindent\textbf{Keywords:} groundwater remediation; chance-constrained
optimization; conditional value-at-risk; physics-informed neural networks;
Fourier neural operators; verification and validation; reproducibility.
\end{abstract}
